\chapter{Interazioni intermolecolari}

Tutto quanto discusso finora riguarda le forze intramolecolari\mkeyword{forze intramolecolari}, cioè i legami covalenti che tengono uniti gli atomi all'interno di una molecola.
Resta però da spiegare perché due molecole distinte, già sature dal punto di vista della valenza, si attraggano comunque.
Se così non fosse nessun gas condenserebbe. Le interazioni responsabili di questa coesione residua si dicono forze intermolecolari\mkeyword{forze intermolecolari}, e sono di natura elettrostatica ma indiretta: nascono tutte dall'accoppiamento fra distribuzioni di carica non uniformi, permanenti o transitorie.\\

L'ordine di grandezza è quello che ci si aspetta da un effetto residuo, cioè qualche $kJ/mol$ contro le centinaia di $kJ/mol$ di un legame covalente.
Questa gerarchia ha una conseguenza operativa immediata: quando una sostanza molecolare fonde o evapora si rompono le interazioni intermolecolari e non i legami covalenti, che restano intatti.\\

Da qui segue anche il criterio di confronto usato in laboratorio.
A parità di peso molecolare, temperature di fusione ed ebollizione più alte segnalano interazioni intermolecolari più intense.
La grandezza fisicamente rilevante è la polarizzabilità\mkeyword{polarizzabilità}, che a sua volta cresce con il numero di elettroni e quindi, indirettamente, con il peso molecolare: la massa in sé non compare in nessuna delle espressioni che ricaveremo.

\section{Interazioni ione-dipolo}

L'interazione più intensa fra le forze intermolecolari è quella ione-dipolo\mkeyword{interazione ione-dipolo}, che si stabilisce fra la carica netta di uno ione e la carica parziale di una molecola polare.

\begin{figure}[htbp]
    \centering
    \begin{tikzpicture}[>={Latex[length=2mm]}]
        % --- macro riutilizzabili -----------------------------------------
        \def\ione#1{%
            \draw (#1) circle (0.45);
            \node at (#1) {$+$};
        }
        \def\dipolo#1#2{%
            \begin{scope}[shift={(#1)}]
                \draw (0,0) ellipse (0.80 and 0.45);
                \draw[->] ({-0.45*#2},0) -- ({0.45*#2},0);
                \node[font=\footnotesize, anchor=south] at ({-0.62*#2},0.52) {$\delta^{+}$};
                \node[font=\footnotesize, anchor=south] at ({0.62*#2},0.52) {$\delta^{-}$};
            \end{scope}
        }
        \def\interazione#1#2{\draw[dotted, thick] (#1,0) -- (#2,0);}
        % --- disegno ------------------------------------------------------
        \ione{0,0}
        \interazione{0.45}{1.80}
        \dipolo{2.60,0}{-1}
    \end{tikzpicture}
    \caption{Interazione ione-dipolo.}
    \label{fig:ione-dipolo}
\end{figure}

L'energia di interazione ha la forma

\begin{equation}
    U_{\text{ione-dipolo}} = -\frac{q\mu\cos\theta}{4\pi\varepsilon_0 r^2},
    \label{eq:ionedipolo}
\end{equation}

dove $\theta$ è l'angolo fra l'asse del dipolo e la congiungente.
Il decadimento come $r^{-2}$, molto più lento di quello di tutte le altre interazioni intermolecolari, spiega da solo perché questa sia la più forte del gruppo e perché il suo raggio d'azione sia il più esteso.

\section{Interazioni fra dipoli permanenti}

Fra due molecole polari e complessivamente neutre si stabilisce l'interazione dipolo-dipolo\mkeyword{interazione dipolo-dipolo}, o forza di Keesom\mkeyword{forza di Keesom}, per cui l'estremità positiva di una molecola tende ad avvicinarsi a quella negativa dell'altra.

\begin{figure}[htbp]
    \centering
    \begin{tikzpicture}[>={Latex[length=2mm]}]
        % --- macro riutilizzabili -----------------------------------------
        \def\dipolo#1#2{%
            \begin{scope}[shift={(#1)}]
                \draw (0,0) ellipse (0.80 and 0.45);
                \draw[->] ({-0.45*#2},0) -- ({0.45*#2},0);
                \node[font=\footnotesize, anchor=south] at ({-0.62*#2},0.52) {$\delta^{+}$};
                \node[font=\footnotesize, anchor=south] at ({0.62*#2},0.52) {$\delta^{-}$};
            \end{scope}
        }
        \def\interazione#1#2{\draw[dotted, thick] (#1,0) -- (#2,0);}
        % --- disegno ------------------------------------------------------
        \dipolo{0,0}{1}
        \interazione{0.80}{2.20}
        \dipolo{3.00,0}{1}
    \end{tikzpicture}
    \caption{Interazione dipolo-dipolo.}
    \label{fig:dipolo-dipolo}
\end{figure}

Nel solido, dove le molecole sono bloccate in un reticolo regolare, l'orientazione è ottimale e l'interazione è massimizzata; nel liquido l'ordine è solo locale e parziale, ma sufficiente a rendere il liquido polare più coeso di un apolare di pari massa.

L'energia di interazione ha la forma 

\begin{equation}
    \langle U \rangle_{\text{Keesom}} = -\frac{2}{3}\,\frac{\mu_1^2\mu_2^2}{(4\pi\varepsilon_0)^2 k_B T}\,\frac{1}{r^6}.
    \label{eq:keesom}
\end{equation}

È l'unico dei contributi intermolecolari a dipendere esplicitamente da $T$, e questa dipendenza è ciò che permette di separarlo sperimentalmente dagli altri misurando il secondo coefficiente del viriale a temperature diverse.

\section{Dipoli indotti e polarizzabilità}

Una molecola immersa in un campo elettrico $\mathbf{E}$ vede la propria nuvola elettronica deformarsi rispetto alla posizione dei nuclei, e acquista di conseguenza un momento di dipolo che, per campi non troppo intensi, è lineare nel campo:

\begin{equation}
    \boldsymbol{\mu}_{\text{ind}} = \alpha\mathbf{E} \qquad \alpha = \left.\frac{\partial\mu}{\partial E}\right|_{E=0}
    \label{eq:polarizzabilita}
\end{equation}

La costante $\alpha$ si dice polarizzabilità\mkeyword{polarizzabilità} e misura quanto la distribuzione elettronica sia cedevole.
La \eqref{eq:polarizzabilita} va letta come il primo termine di uno sviluppo in serie del momento di dipolo nel campo, ed è per questo che la definizione corretta di $\alpha$ è quella come derivata calcolata a campo nullo; i termini successivi, detti iperpolarizzabilità, diventano rilevanti solo nei campi intensi dell'ottica non lineare.

L'andamento periodico di $\alpha$ segue dal fatto che a essere deformata è la parte più esterna e meno trattenuta della nuvola.
Scendendo lungo un gruppo la polarizzabilità cresce, perché gli elettroni di valenza si trovano in gusci sempre più lontani dal nucleo; lungo un periodo diminuisce, perché $Z_{\text{eff}}$ aumenta e il raggio atomico si contrae.

\newpage

Su questo meccanismo si innestano due interazioni.

\subsection{Interazione ione-dipolo indotto}

L'interazione ione-dipolo indotto\mkeyword{ione-dipolo indotto} si ha quando uno ione polarizza una molecola apolare

\begin{figure}[htbp]
    \centering
    \begin{tikzpicture}[>={Latex[length=2mm]}]
        % --- macro riutilizzabili -----------------------------------------
        \def\ione#1{%
            \draw (#1) circle (0.45);
            \node at (#1) {$+$};
        }
        \def\dipoloind#1#2{%
            \begin{scope}[shift={(#1)}]
                \draw[dashed] (0,0) circle (0.62);
                \draw (0,0) ellipse (0.98 and 0.39);
                \draw[->] ({-0.55*#2},0) -- ({0.55*#2},0);
                \node[font=\footnotesize, anchor=south] at ({-0.75*#2},0.66) {$\delta^{+}$};
                \node[font=\footnotesize, anchor=south] at ({0.75*#2},0.66) {$\delta^{-}$};
            \end{scope}
        }
        \def\interazione#1#2{\draw[dotted, thick] (#1,0) -- (#2,0);}
        % --- disegno ------------------------------------------------------
        \ione{0,0}
        \interazione{0.45}{1.77}
        \dipoloind{2.75,0}{-1}
    \end{tikzpicture}
    \caption{Interazione ione-dipolo indotto.
    La circonferenza tratteggiata indica la forma della nuvola elettronica in assenza dello ione.}
    \label{fig:ione-dipolo-indotto}
\end{figure}

Un esempio di questa interazione si ha nel caso di $\mathrm{O_2}$ legato allo ione $Fe^{2+}$ dell'emoglobina.

\subsection{Interazione dipolo-dipolo indotto}

L'interazione dipolo-dipolo indotto\mkeyword{dipolo-dipolo indotto}, o forza di Debye\mkeyword{forza di Debye}, è l'analogo con un dipolo permanente al posto dello ione

\begin{figure}[htbp]
    \centering
    \begin{tikzpicture}[>={Latex[length=2mm]}]
        % --- macro riutilizzabili -----------------------------------------
        \def\dipolo#1#2{%
            \begin{scope}[shift={(#1)}]
                \draw (0,0) ellipse (0.80 and 0.45);
                \draw[->] ({-0.45*#2},0) -- ({0.45*#2},0);
                \node[font=\footnotesize, anchor=south] at ({-0.62*#2},0.52) {$\delta^{+}$};
                \node[font=\footnotesize, anchor=south] at ({0.62*#2},0.52) {$\delta^{-}$};
            \end{scope}
        }
        \def\dipoloind#1#2{%
            \begin{scope}[shift={(#1)}]
                \draw[dashed] (0,0) circle (0.62);
                \draw (0,0) ellipse (0.98 and 0.39);
                \draw[->] ({-0.55*#2},0) -- ({0.55*#2},0);
                \node[font=\footnotesize, anchor=south] at ({-0.75*#2},0.66) {$\delta^{+}$};
                \node[font=\footnotesize, anchor=south] at ({0.75*#2},0.66) {$\delta^{-}$};
            \end{scope}
        }
        \def\interazione#1#2{\draw[dotted, thick] (#1,0) -- (#2,0);}
        % --- disegno ------------------------------------------------------
        \dipolo{0,0}{1}
        \interazione{0.80}{2.17}
        \dipoloind{3.15,0}{1}
    \end{tikzpicture}
    \caption{Interazione dipolo-dipolo indotto}
    \label{fig:dipolo-dipolo-indotto}
\end{figure}

L'energia di interazione ha la forma 

\begin{equation}
    U_{\text{Debye}} = -\frac{\mu_1^2\alpha_2 + \mu_2^2\alpha_1}{(4\pi\varepsilon_0)^2}\,\frac{1}{r^6}.
    \label{eq:debye}
\end{equation}

A differenza della \eqref{eq:keesom} non compare alcuna dipendenza dalla temperatura, perché il dipolo indotto segue istantaneamente l'orientazione del campo e non può essere randomizzato dall'agitazione termica.

\newpage

\section{Forze di dispersione di London}

Il contributo di gran lunga più importante è quello delle forze di dispersione\mkeyword{forze di dispersione}, dette anche forze di London\mkeyword{forze di London}.

Il meccanismo è il seguente: in un dato istante la nuvola elettronica di un atomo, pur essendo mediamente sferica, si trova per fluttuazione spostata da un lato, generando un dipolo istantaneo che polarizza il vicino, e i due dipoli così correlati si attraggono.

\begin{figure}[htbp]
    \centering
    \begin{tikzpicture}[>={Latex[length=2mm]}]
        % --- macro riutilizzabili -----------------------------------------
        \def\apolare#1{\draw (#1) circle (0.62);}
        \def\dipoloist#1#2{%
            \begin{scope}[shift={(#1)}]
                \draw[dashed] (0,0) circle (0.62);
                \draw (0,0) ellipse (0.98 and 0.39);
                \draw[->] ({-0.55*#2},0) -- ({0.55*#2},0);
                \node[font=\footnotesize, anchor=south] at ({-0.75*#2},0.66) {$\delta^{+}$};
                \node[font=\footnotesize, anchor=south] at ({0.75*#2},0.66) {$\delta^{-}$};
            \end{scope}
        }
        \def\interazione#1#2{\draw[dotted, thick] (#1,0) -- (#2,0);}
        % --- pannello (a): nuvole imperturbate ----------------------------
        \apolare{0,0}
        \interazione{0.62}{2.18}
        \apolare{2.80,0}
        \node[font=\footnotesize] at (1.40,-1.55) {(a)};
        % --- pannello (b): fluttuazione e induzione -----------------------
        \dipoloist{6.20,0}{1}
        \interazione{7.18}{8.02}
        \dipoloist{9.00,0}{1}
        \node[font=\footnotesize] at (7.60,-1.55) {(b)};
        \node[font=\footnotesize, anchor=north] at (6.20,-0.80) {istantaneo};
        \node[font=\footnotesize, anchor=north] at (9.00,-0.80) {indotto};
    \end{tikzpicture}
    \caption{Forza di dispersione di London.}
    \label{fig:dispersione}
\end{figure}


Poiché il meccanismo non richiede alcun dipolo permanente, queste forze sono universali e agiscono fra qualunque coppia di sistemi, atomi di gas nobile inclusi.\\

Come nei casi precedenti si ottiene ancora un decadimento come $r^{-6}$ e per due sistemi identici di energia di ionizzazione $I$ si ha

\begin{equation}
    U_{\text{London}} = -\frac{3}{4}\,\frac{I\alpha^2}{(4\pi\varepsilon_0)^2}\,\frac{1}{r^6}.
    \label{eq:london}
\end{equation}

Le forze di London sono le più deboli soltanto nel confronto fra molecole piccole e fortemente polari, dove il termine di Keesom domina; già per $\mathrm{HCl}$ la dispersione contribuisce per circa quattro volte il termine dipolo-dipolo, e per molecole grandi diventa l'unico contributo che conti davvero.

\section{Repulsione a corto raggio e potenziale di Lennard-Jones}

L'attrazione da sola porterebbe al collasso, quindi deve esistere un termine repulsivo che diventa dominante a piccole distanze.
La sua origine è il principio di esclusione di Pauli\mkeyword{repulsione di Pauli}: quando le nuvole elettroniche cominciano a compenetrarsi, gli elettroni di una molecola non possono occupare gli stati già occupati dall'altra, e devono essere promossi a stati di energia superiore.
Il costo energetico è quindi in larga parte un aumento di energia cinetica dovuto all'ortogonalizzazione delle funzioni d'onda, e non un termine coulombiano fra nuclei.\\

La forma fisicamente corretta del termine repulsivo è $V_{\text{rep}} = Ae^{-Br}$, nota come potenziale di Born-Mayer.\\

\newpage

Nella pratica si usa però quasi sempre il potenziale di Lennard-Jones\mkeyword{potenziale di Lennard-Jones}

\begin{equation}
    V(r) = 4\varepsilon\left[\left(\frac{\sigma}{r}\right)^{12} - \left(\frac{\sigma}{r}\right)^{6}\right],
    \label{eq:lennardjones}
\end{equation}

in cui $\varepsilon$ è la profondità della buca, cioè l'energia di legame della coppia, e $\sigma$ la distanza alla quale il potenziale si annulla.

\begin{figure}[htbp]
    \centering
    \begin{tikzpicture}[x=3cm, y=1.3cm, >={Latex[length=2mm]}]
        % --- macro locali -------------------------------------------------
        \def\LJ{4*(pow(1/\x,12)-pow(1/\x,6))}
        \def\rmin{1.1225}
        % --- assi ---------------------------------------------------------
        \draw[->] (0.88,0) -- (3.25,0) node[below right] {$r$};
        \draw[->] (0.92,-1.45) -- (0.92,2.25) node[above left] {$V(r)$};
        % --- curva --------------------------------------------------------
        \draw[domain=0.95:3.15, samples=250, smooth, thick]
            plot (\x, {\LJ});
        % --- sigma --------------------------------------------------------
        \draw[dashed] (1,0) -- (1,-1.15);
        \draw[<->] (0.92,-1.15) -- (1,-1.15);
        \node[above] at (0.96,-1.1) {$\sigma$};
        % --- r_min --------------------------------------------------------
        \draw[dashed] (\rmin,0) -- (\rmin,-1);
        \draw[dashed] (0.92,-1) -- (\rmin,-1);
        \draw[<->] (0.92,-1.35) -- (\rmin,-1.35);
        \node[below] at (1.02,-1.35) {$r_{\text{min}}$};
        % --- epsilon ------------------------------------------------------
        \draw[<->] (1.55,0) -- (1.55,-1);
        \node[right] at (1.55,-0.5) {$\varepsilon$};
        \draw[dashed] (\rmin,-1) -- (1.55,-1);
        % --- etichette dei due regimi -------------------------------------
        \node[right] at (1.05,1.6) {repulsione di Pauli, $\sim r^{-12}$};
        \node[right] at (2.05,0.42) {dispersione, $\sim -r^{-6}$};
        \draw (2.6,0.32) -- (2.6,-0.05);
    \end{tikzpicture}
    \caption{Potenziale di Lennard-Jones.}
    \label{fig:lennardjones}
\end{figure}

Annullando la derivata della \eqref{eq:lennardjones} si trova immediatamente la posizione del minimo,
\begin{equation}
    r_{\text{min}} = 2^{1/6}\sigma \simeq 1{,}12\,\sigma,
    \qquad V(r_{\text{min}}) = -\varepsilon.
    \label{eq:rmin}
\end{equation}

L'esponente $12$ non ha alcun contenuto fisico, ed è importante saperlo.
La repulsione vera è esponenziale, come si è appena detto; il $12$ fu scelto per una ragione esclusivamente computazionale, cioè che $12 = 2\times 6$ e dunque, calcolato una volta il termine $(\sigma/r)^6$, quello repulsivo si ottiene elevandolo al quadrato, senza valutare né esponenziali né radici.
In una simulazione di dinamica molecolare il potenziale va calcolato per ogni coppia di particelle a ogni passo temporale, e questo risparmio faceva la differenza fra fattibile e impraticabile sui calcolatori dell'epoca in cui il modello fu introdotto.
Il $-6$, al contrario, è fisica autentica e discende dalla \eqref{eq:london}: conviene tenere separate le due cose, perché il termine attrattivo è teoria mentre quello repulsivo è convenienza numerica, tanto che i campi di forza moderni più accurati ripristinano la forma esponenziale con il potenziale di Buckingham.

\section{Il legame a idrogeno}
 
Fra tutte le interazioni intermolecolari il legame a idrogeno\mkeyword{legame a idrogeno} è la più complessa.\\

Si presenta di solito come un caso particolare di interazione dipolo-dipolo, e in prima approssimazione la descrizione regge: si stabilisce fra molecole nelle quali un atomo di idrogeno è legato covalentemente a un atomo piccolo e fortemente elettronegativo, in pratica $\mathrm{F}$, $\mathrm{N}$ oppure $\mathrm{O}$, cosicché il legame risulta molto polarizzato e l'idrogeno porta una carica parziale positiva pronunciata.
Vedremo però che questa lettura, per quanto utile, è insufficiente, e che il legame a idrogeno possiede una componente che nessun modello elettrostatico può riprodurre.
 
Si chiama donatore\mkeyword{donatore} il frammento $\mathrm{X}$--$\mathrm{H}$ che mette a disposizione l'idrogeno, e accettore\mkeyword{accettore} l'atomo $\mathrm{Y}$ che offre un doppietto libero, secondo lo schema

\begin{equation}
    \mathrm{X}\text{--}\mathrm{H}\cdots\mathrm{Y},
    \qquad \mathrm{X},\mathrm{Y} \in \{\mathrm{F},\mathrm{N},\mathrm{O}\}.
    \label{eq:schemaH}
\end{equation}

Nell'acqua ogni molecola può fungere da donatore due volte, avendo due idrogeni, e da accettore altre due, avendo due doppietti liberi sull'ossigeno, il che le consente di formare fino a quattro legami a idrogeno contemporaneamente.
Nell'acido fluoridrico, al contrario, il singolo idrogeno limita il numero di legami donati a uno, e nell'ammoniaca è l'unico doppietto libero dell'azoto a limitare a uno il numero di legami accettati.

Questa contabilità è essenziale e spiega un fatto altrimenti paradossale: il singolo legame $\mathrm{F}$--$\mathrm{H}\cdots\mathrm{F}$ è il più forte della famiglia, con circa $29 \ kJ/mol$ contro i $21 \ kJ/mol$ del legame nell'acqua, eppure l'acido fluoridrico bolle a $19,5 \ ^\circ C$ e l'acqua a $100 \ ^\circ C$.
Il motivo è che $\mathrm{HF}$ può costruire soltanto catene monodimensionali a zig-zag, mentre l'acqua costruisce una rete tridimensionale con il doppio dei legami per molecola.
 
\subsection{Perché il legame a idrogeno è così forte}
 
Le ragioni per cui il legame a idrogeno è così forte sono due.\\

La prima è che il legame $\mathrm{X}$--$\mathrm{H}$ è fortemente polarizzato, con una carica parziale $\delta^+$ concentrata sull'idrogeno e $\delta^-$ sull'atomo elettronegativo.

La seconda è che tanto l'idrogeno quanto $\mathrm{N}$, $\mathrm{O}$ e $\mathrm{F}$ sono atomi piccoli, cosicché i due dipoli possono avvicinarsi molto, e poiché l'energia dipolare cresce rapidamente al diminuire della distanza, l'effetto è amplificato.\\

Va aggiunta poi una peculiarità dell'idrogeno, e che è la vera ragione per cui l'interazione porta il suo nome.

L'idrogeno è l'unico elemento privo di elettroni di core: quando cede densità elettronica al partner più elettronegativo, ciò che resta esposto è il nucleo nudo, cioè un protone.
In ogni altro atomo la carica positiva del nucleo è schermata dai gusci interni, e soprattutto quei gusci oppongono la repulsione di Pauli discussa in precedenza, che impedisce all'accettore di avvicinarsi.
Nell'idrogeno questa barriera semplicemente non c'è, e il doppietto dell'accettore può penetrare fino a distanze altrimenti proibite.

\subsection{La componente covalente}
 
La classificazione del legame a idrogeno come caso particolare di interazione dipolo-dipolo va dunque corretta.

Tre osservazioni sperimentali non sono compatibili con un modello puramente elettrostatico.
 
La prima è la direzionalità.
L'energia di un'interazione fra due dipoli varia dolcemente con l'angolo, mentre il legame a idrogeno richiede un allineamento marcato, con l'angolo $\mathrm{X}$--$\mathrm{H}\cdots\mathrm{Y}$ prossimo a $180^\circ$ e un rapido decadimento dell'energia già per deviazioni di qualche decina di gradi.
Un allineamento così stringente è la firma della sovrapposizione fra orbitali, che dipende dall'orientazione reciproca in modo assai più critico di quanto non faccia un'interazione fra cariche.
 
La seconda è l'allungamento del legame del donatore.
Quando si forma il legame a idrogeno, il legame covalente $\mathrm{X}$--$\mathrm{H}$ si allunga e la sua costante di forza diminuisce.
Nessuna interazione puramente elettrostatica ha motivo di indebolire il legame covalente del partner.
 
La terza prova viene dalla risonanza magnetica nucleare, dove la riga di assorbimento di un nucleo si sdoppia se un nucleo vicino gli è accoppiato.
L'accoppiamento può avvenire in due modi: direttamente attraverso lo spazio, per interazione fra i momenti magnetici, oppure per il tramite degli elettroni di legame, e in un liquido il primo si annulla mediando sulle orientazioni mentre il secondo, essendo uno scalare, sopravvive.
Negli acidi nucleici si misura fra i nuclei di azoto di due basi appaiate, separati dal solo legame a idrogeno, un accoppiamento di qualche $Hz$: poiché in soluzione soltanto gli elettroni possono trasmetterlo, ne segue che una densità elettronica condivisa attraversa effettivamente il legame a idrogeno.\\

La descrizione corretta è quella di un legame a tre centri e quattro.
Combinando gli orbitali dei tre centri allineati $\mathrm{X}$, $\mathrm{H}$ e $\mathrm{Y}$ si ottengono un orbitale legante esteso a tutti e tre, uno all'incirca non legante e uno antilegante; i quattro elettroni disponibili, due del legame $\sigma_{\mathrm{X-H}}$ e due del doppietto di $\mathrm{Y}$, riempiono i primi due e lasciano vuoto il terzo, con un guadagno netto di energia.
Equivalentemente, il doppietto di $\mathrm{Y}$ si delocalizza in parte nell'orbitale $\sigma^*$ del legame $\mathrm{X}$--$\mathrm{H}$: poiché quell'orbitale ha un lobo sull'idrogeno e un nodo fra $\mathrm{X}$ e $\mathrm{H}$, la sovrapposizione è legante attraverso $\mathrm{H}\cdots\mathrm{Y}$ e antilegante attraverso $\mathrm{X}$--$\mathrm{H}$.
Una sola delocalizzazione produce dunque due effetti opposti, ed è questa la ragione per cui il legame a idrogeno che si forma e l'allungamento del legame covalente del donatore sono la stessa cosa vista da due lati.
Ne segue una predizione verificata: al crescere della forza del legame a idrogeno la distanza $\mathrm{H}\cdots\mathrm{Y}$ si accorcia, quella $\mathrm{X}$--$\mathrm{H}$ si allunga.
Resta inteso che il contributo elettrostatico rimane maggioritario nei legami di forza ordinaria, dove si stima intorno al $60$--$80\%$: la componente covalente è minoritaria, ma è quella che rende conto di tutto ciò che l'elettrostatica non spiega.

\subsection{Il ghiaccio}
 
Di norma un solido è più denso del proprio liquido, perché la perdita di energia cinetica consente un impacchettamento più efficiente.
L'acqua è la nota eccezione, con $0{,}917 \ g/cm^3$ per il ghiaccio contro $1{,}000$ per il liquido a $0 \ ^\circ C$, ed è il legame a idrogeno a produrre l'inversione.
 
Nel ghiaccio ordinario ogni molecola d'acqua forma esattamente quattro legami a idrogeno, due donati e due accettati, diretti verso i vertici di un tetraedro: ciascun idrogeno punta verso un doppietto libero di una molecola vicina, e ciascun doppietto riceve un idrogeno.
La struttura che ne risulta è un reticolo di anelli esagonali con un'ampia cavità centrale, tanto ampia che una molecola d'acqua potrebbe esservi inscritta.
La chiave del fenomeno sta nel fatto che questa disposizione ottimizza il numero e la geometria dei legami a idrogeno ma non l'impacchettamento, e le due cose sono in conflitto: l'angolo tetraedrico di $109{,}5^\circ$ imposto dai doppietti dell'ossigeno obbliga le molecole a tenersi distanti fra loro, e il volume del solido ne risulta penalizzato.
Fondendo, una frazione di questi legami si rompe, il reticolo collassa parzialmente e le molecole possono avvicinarsi, con un aumento netto di densità.
Il processo prosegue anche sopra il punto di fusione, ed è per questo che l'acqua liquida raggiunge il massimo di densità a $4 \ ^\circ C$ e non a $0$: fra $0$ e $4$ gradi la progressiva demolizione dei residui di struttura tetraedrica prevale sulla dilatazione termica ordinaria, che sopra i $4$ gradi torna a dominare.

\section{Quadro riassuntivo}

La Tabella~\ref{tab:riepilogo} raccoglie gli ordini di grandezza tipici, che vanno però letti con la cautela discussa: gli intervalli si sovrappongono ampiamente, e l'importanza relativa di ciascun contributo in un sistema concreto dipende dalle dimensioni delle molecole coinvolte assai più che dalla classificazione formale.

\begin{table}[htbp]
    \centering
    \begin{tabular}{llcl}
        \toprule
        Interazione & Grandezze coinvolte & $E \ [kJ/mol]$ & Dipendenza \\
        \midrule
        Ione-dipolo          & carica e dipolo permanente     & $40$--$600$   & $r^{-2}$ \\
        Legame a idrogeno    & $\mathrm{X}$--$\mathrm{H}$ e doppietto & $10$--$40$ & direzionale \\
        Dipolo-dipolo        & due dipoli permanenti          & $5$--$25$     & $r^{-6}$, $\propto T^{-1}$ \\
        Ione-dipolo indotto  & carica e polarizzabilità       & $3$--$15$     & $r^{-4}$ \\
        Dipolo-dipolo indotto& dipolo e polarizzabilità       & $2$--$10$     & $r^{-6}$ \\
        Dispersione          & due polarizzabilità            & $0{,}05$--$40$& $r^{-6}$ \\
        \bottomrule
    \end{tabular}
    \caption{Ordini di grandezza delle interazioni intermolecolari.
    L'intervallo estremamente ampio della dispersione riflette il fatto che essa cresce senza limite con le dimensioni molecolari.}
    \label{tab:riepilogo}
\end{table}
